
\documentclass{article}
\usepackage{colortbl}
\usepackage{makecell}
\usepackage{multirow}
\usepackage{supertabular}

\begin{document}

\newcounter{utterance}

\centering \large Interaction Transcript for game `codenames', experiment `abstract', episode 8 with reflexion{-}Qwen3.5{-}9B.
\vspace{24pt}

{ \footnotesize  \setcounter{utterance}{1}
\setlength{\tabcolsep}{0pt}
\begin{supertabular}{c@{$\;$}|p{.15\linewidth}@{}p{.15\linewidth}p{.15\linewidth}p{.15\linewidth}p{.15\linewidth}p{.15\linewidth}}
    \# & $\;$A & \multicolumn{4}{c}{Game Master} & $\;\:$B\\
    \hline

    \theutterance \stepcounter{utterance}  
    & & & \multicolumn{2}{p{0.3\linewidth}}{
        \cellcolor[rgb]{0.9,0.9,0.9}{
            \makecell[{{p{\linewidth}}}]{
                \texttt{\tiny{[GM$|$GM]}}
                {'hidden': {'team': ['dilution', 'existentialism'], 'innocent': ['stigma', 'ridiculousness', 'optimism', 'inefficiency', 'exclusivity', 'adorability'], 'opponent': [], 'assassin': ['coercion']}, 'revealed': {'team': {'team': [], 'innocent': [], 'opponent': [], 'assassin': []}, 'opponent': {'team': [], 'innocent': [], 'opponent': [], 'assassin': []}}}
            }
        }
    }
    & & \\ \\

    \theutterance \stepcounter{utterance}  
    & & \multicolumn{4}{p{0.6\linewidth}}{
        \cellcolor[rgb]{0.9,0.9,0.9}{
            \makecell[{{p{\linewidth}}}]{
                \texttt{\tiny{[P1$\langle$GM]}}
                \texttt{You have attempted this same task before and did not succeed. These are the plans you wrote after those attempts.} \\
\texttt{Trial \#1: New plan: I failed because I included "mix" as a clue word, which is explicitly listed as an opponent's word. The rules state the clue cannot be one of the words in the lists. Additionally, "mix" is semantically too broad and directly overlaps with the opponent's word "intermix," risking a guess of an opponent word. Next time, I will strictly verify that my chosen clue word is not present in any of the provided lists (team, opponent, distractor, or kill words) and ensure it is a distinct semantic bridge rather than a direct synonym or component of the target words. I will also double{-}check the format to ensure no extra text is included.} \\
\\ 
\texttt{Trial \#2: New plan: I failed because my clue "thin" was too specific to "dilution" and failed to connect semantically to "existentialism," causing my partner to guess a distractor word ("unsettlement") instead of the remaining target. This wasted a turn and allowed the opponent to guess "fruition." Next time, I will choose a broader, high{-}level concept that encompasses all remaining targets simultaneously (e.g., "philosophy" for existentialism and dilution, or "state" for both) to maximize the chance of clearing multiple words in one turn without risking a distractor guess. I will also ensure the clue is distinct from all opponent and distractor words to avoid triggering an innocent guess that ends the turn prematurely.} \\
\\ 
\\ 
\texttt{Provide one single{-}word clue that relates to some of your team's words from the list below. You can choose to how many and to which words your clue relates to.} \\
\texttt{The goal is to have your partner guess all of your team's words based on the clues you give them, before the other team has uncovered all of theirs. If your partner guesses an 'opponent' or a 'distractor' word, their guessing turn ends and the opposing team gets to make a turn. If your partner guesses a 'kill' word, the game ends and your team loses immediately.} \\
\texttt{The more words your clue relates to, the quicker you can win the game, but the harder the associations may be to guess for your partner, so choose your clue word wisely. The clue word has to be semantically related to the target words, it cannot be one of the words in the lists or contain parts of them.} \\
\\ 
\texttt{Always give your single{-}word clue and your comma{-}separated list of related target words in the following format and make your answers as short as possible, never include any other text than is required in this form:} \\
\\ 
\texttt{CLUE: <WORD>} \\
\texttt{TARGETS: <TARGETS>} \\
\\ 
\texttt{Your team words are: dilution, existentialism.} \\
\texttt{Your opponent's team words are: .} \\
\texttt{Distractor words are: stigma, ridiculousness, optimism, inefficiency, exclusivity, adorability.} \\
\texttt{The kill words are: coercion.} \\
            }
        }
    }
    & \\ \\

    \theutterance \stepcounter{utterance}  
    & \multicolumn{4}{p{0.6\linewidth}}{
        \cellcolor[rgb]{0.9,0.9,0.9}{
            \makecell[{{p{\linewidth}}}]{
                \texttt{\tiny{[P1$\rangle$GM]}}
                \texttt{CLUE: philosophy} \\
\texttt{TARGETS: existentialism,dilution} \\
            }
        }
    }
    & & \\ \\

    \theutterance \stepcounter{utterance}  
    & & & \multicolumn{2}{p{0.3\linewidth}}{
        \cellcolor[rgb]{0.9,0.9,0.9}{
            \makecell[{{p{\linewidth}}}]{
                \texttt{\tiny{[GM$|$GM]}}
                \texttt{philosophy} \\
            }
        }
    }
    & & \\ \\

    \theutterance \stepcounter{utterance}  
    & & & \multicolumn{2}{p{0.3\linewidth}}{
        \cellcolor[rgb]{0.9,0.9,0.9}{
            \makecell[{{p{\linewidth}}}]{
                \texttt{\tiny{[GM$|$GM]}}
                ['existentialism', 'dilution']
            }
        }
    }
    & & \\ \\

    \theutterance \stepcounter{utterance}  
    & & & \multicolumn{2}{p{0.3\linewidth}}{
        \cellcolor[rgb]{0.9,0.9,0.9}{
            \makecell[{{p{\linewidth}}}]{
                \texttt{\tiny{[GM$|$GM]}}
                {'word': 'existentialism', 'assignment': 'team'}
            }
        }
    }
    & & \\ \\

    \theutterance \stepcounter{utterance}  
    & & & \multicolumn{2}{p{0.3\linewidth}}{
        \cellcolor[rgb]{0.9,0.9,0.9}{
            \makecell[{{p{\linewidth}}}]{
                \texttt{\tiny{[GM$|$GM]}}
                {'word': 'dilution', 'assignment': 'team'}
            }
        }
    }
    & & \\ \\

    \theutterance \stepcounter{utterance}  
    & & & \multicolumn{2}{p{0.3\linewidth}}{
        \cellcolor[rgb]{0.9,0.9,0.9}{
            \makecell[{{p{\linewidth}}}]{
                \texttt{\tiny{[GM$|$GM]}}
                \texttt{opponent has won} \\
            }
        }
    }
    & & \\ \\

    \theutterance \stepcounter{utterance}  
    & \multicolumn{4}{p{0.6\linewidth}}{
        \cellcolor[rgb]{0.95,0.95,0.95}{
            \makecell[{{p{\linewidth}}}]{
                \texttt{\tiny{[P2$\rightarrow$P1]}}
                \texttt{Post{-}episode Reflexion (not sent to GM):} \\
\texttt{New plan: I failed because "philosophy" is too abstract and likely triggered a distractor guess (e.g., "existentialism" is a branch of philosophy, but "dilution" is not, causing confusion or a wrong guess). Additionally, "philosophy" might be too broad, leading my partner to guess a distractor word like "optimism" or "ridiculousness" instead of the specific targets. Next time, I will choose a more concrete, shared semantic bridge that directly links both "dilution" and "existentialism" without being a category they both belong to. For example, I could use "concentration" (the opposite of dilution, and existentialism deals with the concentration of meaning/existence) or "presence" (dilution reduces presence; existentialism focuses on being/presence). I will also ensure the clue is not a distractor word itself and verify it doesn't accidentally match any opponent or kill words. Finally, I will keep the clue short and specific to avoid ambiguity.} \\
            }
        }
    }
    & & \\ \\

\end{supertabular}
}

\end{document}
